\section{Conclusion and discussion}

We developed an NTK analysis for low-rank random-feature architectures (RF-LR) under explicit assumptions, giving an explicit new recursion, sharp depth scaling for the proxy kernel, condition-number bounds, and---in the three-layer case---RKHS equivalence with the shallow ReLU kernel and concentration in the bottleneck dimension. For kernel regression in the lazy/NTK regime with mean-zero targets, stability and convergence depend on the condition number \(\lambda_{\max}/\lambda_{\min}\) of the centered Gram matrix. Our results give a lower bound on the proxy condition number and, in one setting, its exact value with high probability; we do not claim to govern optimization, and the link to GD convergence or sample complexity is left for future work.

\paragraph{Open directions.}
On the theoretical side: extending concentration bounds to general depth \(L\) requires controlling nested conditional means and \(O(L^2)\) cross-terms from the recursion; a non-Gaussian propagation analysis would give complementary insights. Establishing bulk spectral laws for the empirical RF-LR Gram matrix in the RMT limit (\(n\sim N^2\)), adapting~\cite{benigni2025eigenvaluedistributionneuraltangent} to low-rank depth \(L\ge 2\), is another direction. Conditioning results hold for arbitrary \(L\); RKHS equivalence is proved only for three layers---extending it to \(L\ge 4\) is open (Appendix~\ref{appendix:what-extends-L}). Other extensions include RKHS identification at arbitrary depth, EOC parameterization and \(Nr\) scaling, and data geometry (hypercube vs.\ spherical).

Experimentally, next steps include confirming lazy-kernel predictions and the boundary of lazy training; width sweeps to test \(O(1/n)+O(1/r)\) finite-width corrections and minimal \((n,r)\) for kernel predictions; and comparison with MLPs on conditioning and optimization time at matched parameter counts.

Our scope is the lazy/NTK regime, where tractability is greatest. Completing finite-width formulas, spectrum characterizations, and the experimental program above would strengthen the link between the asymptotic theory and practice.
